\subsection{INVARIANT ELEMENTS}

DEFINITION 3. Let \(\mathfrak{g}\) be a Lie algebra and M a \(\mathfrak{g}\) -module. An element \(m \in \mathbf{M}\) is called invariant (with respect to the \(\mathfrak{g}\) -module structure on M or with respect to the corresponding representation of \(\mathfrak{g}\) ) if \(x_{\mathbf{M}}.m = 0\) for all \(x \in \mathfrak{g}\) .

*Let G be a connected real Lie group, g its Lie algebra, \(\theta\) an analytic representation of G on a finite-dimensional real vector space E and \(\rho\) the corresponding representation of g on E. Let \(m \in E\) . The element m is invariant with respect to \(\rho\) if and only if \(\theta(g).m = m\) for all \(g \in G\) . This justifies the use of the word ``invariant''.*

Example 1. Let M, N be two g-modules and \(P = \mathcal{L}_{K}(M, N)\) . For an element f of P to be invariant, it is necessary and sufficient, by (6), that f be a homomorphism of the g-module M into the g-module N. In particular, if M = N and \(x_{M} = x_{N}\) for all \(x \in g\) , f is invariant if and only if f is permutable with the \(x_{M}\) .

Example 2. Let M be a K-module with a finite basis. If M has a g-module structure, \(\mathcal{L}(\mathrm{M},\mathrm{M})\) and \(\mathbf{M}^*\otimes \mathbf{M}\) have g-module structures and the canonical mapping of \(\mathbf{M}^{*}\otimes \mathbf{M}\) into \(\mathcal{L}(\mathrm{M},\mathrm{M})\) is a g-module isomorphism (Proposition 4). As \(1\in \mathcal{L}(\mathrm{M},\mathrm{M})\) is obviously an invariant (cf.~Example 1), the corresponding element \(u\) of \(\mathbf{M}^{*}\otimes \mathbf{M}\) is an invariant. If \((e_i)_{1\leq i\leq n}\) is a basis of M and \((e_{*}^{i})_{1\leqslant i\leqslant n}\) is the dual basis, we have \(u = \sum_{i = 1}^{n}e_{i}^{*}\otimes e_{i}\) .

Example 3. Let M be a g-module. Let \(\beta\) be a bilinear form on M and \(f\) the corresponding element of \(\mathcal{L}(\mathrm{M},\mathrm{M}^*)\) . For \(\beta\) to be invariant, it is necessary and sufficient that \(f\) be a g-module homomorphism (Proposition 4 and Example 1). Suppose that K is a field and that \(\dim_{\mathbb{K}}\mathrm{M} < +\infty\) . A non-degenerate invariant bilinear form \(\beta\) on M defines an isomorphism of the g-module M onto the g-module \(\mathrm{M}^*\) and hence an isomorphism of the g-module \(\mathrm{M}\otimes \mathrm{M}\) onto the g-module \(\mathrm{M}^*\otimes \mathrm{M}\) . Thus, by Example 2, giving \(\beta\) defines canonically an invariant element \(c\) in the g-module \(\mathrm{M}\otimes \mathrm{M}\) , which can be constructed as follows: let \((e_i)_{1\leqslant i\leqslant n}\) be a basis of M and \((e_i')_{1\leqslant i\leqslant n}\) the basis of M such that \(\beta (e_i,e_j') = \delta_{i,j}\) ; then \(c = \sum_{i=1}^{n}e_i\otimes e_i'\) .

PROPOSITION 5. Let g be a Lie K-algebra, h an ideal of g, \(\rho\) a representation of g on M and \(\rho'\) the restriction of \(\rho\) to h. Then the set N of elements of M invariant with respect to \(\rho'\) is stable under \(\rho(\mathfrak{g})\) .

Let \(n\in \mathbf{N}\) and \(y\in \mathfrak{g}\) ; for all \(x\in \mathfrak{h}\) , \([x,y]\in \mathfrak{h}\) and hence

\[
\rho (x) \rho (y) n = \rho ([ x, y ]) n + \rho (y) \rho (x) n = 0;
\]

hence \(\rho(y)n\in\mathbf{N}\) .

PROPOSITION 6. Let M be a semi-simple g-module. Then the submodule \(M_{0}\) of invariant elements of M admits one and only one supplement stable under the \(x_{M}\) , namely the submodule \(M_{1}\) generated by the \(x_{M}.m\) ( \(x \in g, m \in M\) ).

Let \(\mathbf{M}'\) be a submodule of \(\mathbf{M}\) which is stable under the \(x_{\mathrm{M}}\) and a supplement of \(\mathrm{M}_0\) in \(\mathbf{M}\) . For all \(m \in \mathbf{M}\) , \(m = m_0 + m'\) with \(m_0 \in \mathrm{M}_0\) , \(m' \in \mathrm{M}'\) , and hence \(x_{\mathrm{M}}m = x_{\mathrm{M}}m' \in \mathrm{M}'\) . Hence \(\mathrm{M}_1 \subset \mathrm{M}'\) . Let \(\mathrm{M}_2\) be a submodule of \(\mathbf{M}'\) stable under the \(x_{\mathrm{M}}\) and supplementary to \(\mathrm{M}_1\) in \(\mathrm{M}'\) . For all \(m \in \mathrm{M}_2\) ,

\[
x _ {\mathrm{M}} m \in \mathrm{M} _ {2} \cap \mathrm{M} _ {1} = \{0 \}
\]

for all \(x \in \mathfrak{g}\) , hence \(m \in \mathbf{M}_0\) and hence \(m = 0\) . Hence \(\mathbf{M}_2 = \{0\}\) , which proves that \(\mathbf{M}_1 = \mathbf{M}'\) .

